% GENERATED FILE. Edit canonical lessons, then run npm run generate:books.
% canonical-source-hash: fnv1a64:63cf521fce0a2c0b
% canonical-lessons: ML-C262-kurippu, ML-C262-udaharanam, ML-C262-kalakaran, ML-C262-pujari, ML-C262-gayakan

\chapter{Note, Example, Artist, Temple priest, Singer}
\label{ch:ml-note-example-artist-temple-priest-singer}
\hlchaptermodality{262}

\begin{chapteropening}
\textbf{By the end of this chapter:} \emph{I can say a note, an example, an artist, a temple priest and a singer.}

Complete the last lesson of chapter 262: I can say a note, an example, an artist, a temple priest and a singer.
\end{chapteropening}

\section[\texorpdfstring{\ml{കുറിപ്പ്} (kuṟippŭ)}{കുറിപ്പ് (kuṟippŭ)}]{\ml{കുറിപ്പ്} (kuṟippŭ) — a note}

\label{lesson:ML-C262-kurippu}



\begin{quote}
Before the new one: say the Malayalam for training, then the Malayalam for a report.
\end{quote}

\subsection*{You'll want to know: \ml{കുറിപ്പ്}}
\textbf{\ml{കുറിപ്പ്}} — \emph{kuṟippŭ} — \textquotedblleft{}a note\textquotedblright{}.

\begin{grammarlens}[title={What you've built}]
One more word for this chapter.
\end{grammarlens}

\subsection*{Your turn}
\begin{itemize}
\raggedright
  \item \emph{Say it:} \emph{kuṟippŭ}
  \item \emph{Say it:} \emph{kuṟippŭ}, once more
  \item \emph{Say it:} say \emph{kuṟippŭ}
  \item \emph{Recall:} say the Malayalam for training, then the Malayalam for a report, then say \emph{kuṟippŭ} again
\end{itemize}

\begin{tcolorbox}[breakable,colback=teal!4,colframe=teal!35!black,title={Before you move on}]
What does \ml{കുറിപ്പ്} mean? (\textquotedblleft{}A note\textquotedblright{}.) Say it once more.
\end{tcolorbox}
\section[\texorpdfstring{\ml{ഉദാഹരണം} (udāharaṇaṁ)}{ഉദാഹരണം (udāharaṇaṁ)}]{\ml{ഉദാഹരണം} (udāharaṇaṁ) — an example}

\label{lesson:ML-C262-udaharanam}



\begin{quote}
Before the new one: say the Malayalam for a report, then the Malayalam for a note.
\end{quote}

\subsection*{You'll want to know: \ml{ഉദാഹരണം}}
\textbf{\ml{ഉദാഹരണം}} — \emph{udāharaṇaṁ} — \textquotedblleft{}an example\textquotedblright{}.

\begin{grammarlens}[title={What you've built}]
One more word for this chapter.
\end{grammarlens}

\subsection*{Your turn}
\begin{itemize}
\raggedright
  \item \emph{Say it:} \emph{udāharaṇaṁ}
  \item \emph{Say it:} \emph{udāharaṇaṁ}, once more
  \item \emph{Say it:} say \emph{udāharaṇaṁ}
  \item \emph{Recall:} say the Malayalam for a report, then the Malayalam for a note, then say \emph{udāharaṇaṁ} again
\end{itemize}

\begin{tcolorbox}[breakable,colback=teal!4,colframe=teal!35!black,title={Before you move on}]
What does \ml{ഉദാഹരണം} mean? (\textquotedblleft{}An example\textquotedblright{}.) Say it once more.
\end{tcolorbox}
\section[\texorpdfstring{\ml{കലാകാരൻ} (kalākāran)}{കലാകാരൻ (kalākāran)}]{\ml{കലാകാരൻ} (kalākāran) — an artist}

\label{lesson:ML-C262-kalakaran}



\begin{quote}
Before the new one: say the Malayalam for a note, then the Malayalam for an example.
\end{quote}

\subsection*{You'll want to know: \ml{കലാകാരൻ}}
\textbf{\ml{കലാകാരൻ}} — \emph{kalākāran} — \textquotedblleft{}an artist\textquotedblright{}.

\begin{grammarlens}[title={What you've built}]
One more word for this chapter.
\end{grammarlens}

\subsection*{Your turn}
\begin{itemize}
\raggedright
  \item \emph{Say it:} \emph{kalākāran}
  \item \emph{Say it:} \emph{kalākāran}, once more
  \item \emph{Say it:} say \emph{kalākāran}
  \item \emph{Recall:} say the Malayalam for a note, then the Malayalam for an example, then say \emph{kalākāran} again
\end{itemize}

\begin{tcolorbox}[breakable,colback=teal!4,colframe=teal!35!black,title={Before you move on}]
What does \ml{കലാകാരൻ} mean? (\textquotedblleft{}An artist\textquotedblright{}.) Say it once more.
\end{tcolorbox}
\section[\texorpdfstring{\ml{പൂജാരി} (pūjāri)}{പൂജാരി (pūjāri)}]{\ml{പൂജാരി} (pūjāri) — a temple priest}

\label{lesson:ML-C262-pujari}



\begin{quote}
Before the new one: say the Malayalam for an example, then the Malayalam for an artist.
\end{quote}

\subsection*{You'll want to know: \ml{പൂജാരി}}
\textbf{\ml{പൂജാരി}} — \emph{pūjāri} — \textquotedblleft{}a temple priest\textquotedblright{}.

\begin{grammarlens}[title={What you've built}]
One more word for this chapter.
\end{grammarlens}

\subsection*{Your turn}
\begin{itemize}
\raggedright
  \item \emph{Say it:} \emph{pūjāri}
  \item \emph{Say it:} \emph{pūjāri}, once more
  \item \emph{Say it:} say \emph{pūjāri}
  \item \emph{Recall:} say the Malayalam for an example, then the Malayalam for an artist, then say \emph{pūjāri} again
\end{itemize}

\begin{tcolorbox}[breakable,colback=teal!4,colframe=teal!35!black,title={Before you move on}]
What does \ml{പൂജാരി} mean? (\textquotedblleft{}A temple priest\textquotedblright{}.) Say it once more.
\end{tcolorbox}
\section[\texorpdfstring{\ml{ഗായകൻ} (gāyakan)}{ഗായകൻ (gāyakan)}]{\ml{ഗായകൻ} (gāyakan) — a singer}

\label{lesson:ML-C262-gayakan}



\begin{quote}
Before the new one: say the Malayalam for an artist, then the Malayalam for a temple priest.
\end{quote}

\subsection*{You'll want to know: \ml{ഗായകൻ}}
\textbf{\ml{ഗായകൻ}} — \emph{gāyakan} — \textquotedblleft{}a singer\textquotedblright{}.

\begin{grammarlens}[title={What you've built}]
One more word for this chapter.
\end{grammarlens}

\subsection*{Your turn}
\begin{itemize}
\raggedright
  \item \emph{Say it:} \emph{gāyakan}
  \item \emph{Say it:} \emph{gāyakan}, once more
  \item \emph{Say it:} say \emph{gāyakan}
  \item \emph{Recall:} say the Malayalam for an artist, then the Malayalam for a temple priest, then say \emph{gāyakan} again
\end{itemize}

\begin{tcolorbox}[breakable,colback=teal!4,colframe=teal!35!black,title={Before you move on}]
What does \ml{ഗായകൻ} mean? (\textquotedblleft{}A singer\textquotedblright{}.) Say it once more.
\end{tcolorbox}
